% Integration fragment: uses the main manuscript's proposition/proof environments.
% Place after deterministic adaptive recovery; replace its short randomized example.
% Requires bibliography keys kawase2024 and mastin2015 (see bibitem_suggestions.tex).
\section{Randomized Recovery under a Hard Budget}
\label{sec:randomized-recovery}

Randomization permits a different guarantee from deterministic minimax
recovery. We retain exact disjoint totals, nonnegative uncapped targets,
unit query costs, and one fixed forecast contrast. All information and
deterministic deductions available before recovery are incorporated first.
Remove zero-energy, singleton, and constant-coefficient blocks, whose width
is already zero. Index the remaining blocks by $g=1,\ldots,G$.
If none remain, the optimal width is zero and no query is needed.
Let $\mathcal Q_B$ contain batches with at most $B$ queries in total and
at most $m_g-1$ in block $g$; querying its final coordinate is redundant.
A policy selects $Q\sim\pi$ before reading its answers. The hidden record
$y\in\F$ is fixed before this random draw: the adverse completion may
depend on $\pi$, but cannot depend on the realized batch.

\subsection{Expected width and an exact finite linear program}
Write $S_g(Q)=\M_g\setminus Q$, let $d_g(S)$ denote the coefficient range
on $S$, and set $d_g(S)=0$ for a singleton or empty set. Define
$v_{gi}(Q)=\boldsymbol{1}\{i\in S_g(Q)\}d_g(S_g(Q))$.
Equation~\eqref{eq:querywidth} gives
\begin{equation}
 W(Q,y)=\sum_g\sum_{i\in\M_g}y_i v_{gi}(Q),\qquad
 R_B=\min_{\pi\in\Delta(\mathcal Q_B)}\sup_{y\in\F}
       \mathbb E_{Q\sim\pi}W(Q,y).
 \label{eq:random-objective}
\end{equation}
Here $\Delta(\mathcal Q_B)$ is the set of distributions over admissible
batches. The budget holds for every draw, rather than only in expectation.
The coefficients $c_i$ already include the MSE normalization in
\eqref{eq:linear}.

\begin{proposition}[Finite characterization]
\label{prop:random-finite}
For every batch distribution,
\begin{equation}
 \sup_{y\in\F}\mathbb E_\pi W(Q,y)
 =\sum_g E_g\max_{i\in\M_g}\mathbb E_\pi v_{gi}(Q).
 \label{eq:random-value}
\end{equation}
Consequently $R_B$ is the value of the primal and dual programs
\begin{equation}
 \begin{aligned}
 \min_{\pi,t}\quad &\sum_g E_g t_g,\\
 \text{subject to}\quad
 &\sum_Q\pi_Qv_{gi}(Q)\leq t_g\quad(g,i),\qquad
   \sum_Q\pi_Q=1,\quad\pi_Q\geq0;
 \end{aligned}
 \label{eq:random-primal}
\end{equation}
\begin{equation}
 \begin{aligned}
 \max_{w,\eta}\quad &\eta,\\
 \text{subject to}\quad
 &\eta\leq\sum_g\sum_i w_{gi}v_{gi}(Q)\quad(Q\in\mathcal Q_B),\\
 &\sum_i w_{gi}=E_g,\quad w_{gi}\geq0.
 \end{aligned}
 \label{eq:random-dual}
\end{equation}
The variables $t_g$ and $\eta$ are free; the primal inequalities imply
$t_g\geq0$.
\end{proposition}
\begin{proof}
Interchanging finite sums makes the expected width linear in each
block's target vector. Its maximum over a nonnegative simplex of mass
$E_g$ places that mass on a largest expected coefficient. The blocks are
independent, proving \eqref{eq:random-value}. Introducing their maxima
gives the primal. Linear-programming duality gives the stated dual;
the primal is feasible and has a finite optimum.
\end{proof}

Mixed-strategy robust optimization and dual separation are established
methods. Kawase and Sumita provide a general framework based on weighted
scenario optimization \cite{kawase2024}; Mastin, Jaillet, and Chin study
randomized minmax regret and improvements over deterministic strategies
\cite{mastin2015}. The specialization here is the conditional-width
geometry, the exact coordinate-recovery oracle and compact formulation
below, and the sharp comparison with deterministic adaptive recovery.

\subsection{An exact pricing oracle}
To separate \eqref{eq:random-dual}, minimize
$\sum_{g,i}w_{gi}v_{gi}(Q)$ over admissible batches. Within a block, order
coefficients increasingly, breaking ties by original coordinate index.
For local budget $k$, put $r=m_g-k$. If $r\geq2$, fix the first and last
retained positions $l<u$, with $u-l+1\geq r$. The least cost for these
endpoints is
\begin{equation}
 (c_{g,(u)}-c_{g,(l)})
 \left(w_{g,(l)}+w_{g,(u)}+
       \sum_{i\in T_{g,l,u,r}}w_{g,(i)}\right),
 \label{eq:random-pricing}
\end{equation}
where $T_{g,l,u,r}$ contains the $r-2$ smallest interior weights.
Every retained subset has such endpoints; once they are fixed, its range
is fixed and these interior choices minimize its nonnegative weighted
cost. Enumerate the endpoints, assign cost zero when $r=1$, and apply
the budget-allocation recurrence across blocks. This supplies exact
polynomial-time separation for rational input. Unlike deterministic
minimum-range recovery, an optimal weighted retained set need not be
a consecutive coefficient window.

\subsection{A compact program and an executable batch distribution}
The following formulation avoids a variable for every possible batch.
Set $q=\min(B,\sum_g(m_g-1))$. Any smaller batch can be padded to $q$
queries without widening its conditional interval for any record, so
requiring exactly $q$ loses no optimality for this objective.
The case $B=0$ is obtained directly from the original width formula.

Construct a directed acyclic graph with nodes $(g,b)$ for completed
blocks $0\leq g\leq G$ and budgets $0\leq b\leq q$. For each admissible
$k$ and endpoint pair $(l,u)$ of a retained set of size $r=m_g-k\geq2$,
create an edge $e=(g,b,k,l,u)$ from $(g-1,b)$ to $(g,b+k)$.
For $r=1$, use a single canonical zero-width edge. Impose nonnegative
unit flow $p_e$ from $(0,0)$ to $(G,q)$, with conservation at all other
nodes. Impossible nodes and edges may be omitted.

For an endpoint edge, introduce interior joint inclusion masses satisfying
\begin{equation}
 0\leq z_{ei}\leq p_e\quad(l<i<u),\qquad
 \sum_{l<i<u}z_{ei}=(r-2)p_e.
 \label{eq:random-hypersimplex}
\end{equation}
Its endpoint inclusion masses are $p_e$; interior inclusion masses are
$z_{ei}$, and all others are zero. Denote these by $\alpha_{ei}$ and
put $d_e=c_{g,(u)}-c_{g,(l)}$. Zero-width edges contribute zero.
The compact program minimizes $\sum_gE_gt_g$ subject to the flow
constraints, \eqref{eq:random-hypersimplex}, and
\begin{equation}
 \sum_{e\text{ for block }g}d_e\alpha_{ei}\leq t_g
 \quad(g,i).
 \label{eq:random-compact}
\end{equation}

\begin{proposition}[Exact compact representation]
\label{prop:random-compact}
This program has value $R_B$ and yields an attaining distribution with
at most $B$ queries in every batch. For $B\geq1$, a safe unpruned bound
on its numbers of variables and constraints is
$O(B\sum_g\min(B,m_g)m_g^3)$.
For one block with a fixed local budget, only $O(m_g^3)$ variables are needed.
\end{proposition}
\begin{proof}
A distribution over batches induces flows along their budget paths and
the corresponding endpoint and interior inclusion masses, satisfying
all constraints. Conversely, a nonnegative unit flow on an acyclic graph
decomposes into source-to-sink paths. For any edge with $p_e>0$,
the vector $(z_{ei}/p_e)_{l<i<u}$ belongs to
$\{x\in[0,1]^{u-l-1}:\sum_i x_i=r-2\}$. This hypersimplex is the convex
hull of the indicators of its $(r-2)$-element subsets. To see integrality,
any nonintegral feasible vector has at least two fractional entries;
opposite small perturbations of those entries express it as the midpoint
of two distinct feasible vectors, so it is not a vertex.
Thus each positive-flow edge admits a distribution of legal interiors.
Sample a flow path and, conditionally on each selected edge, such an
interior subset. Every resulting batch uses exactly $q$ queries and
has the coordinate coefficients in \eqref{eq:random-compact}. Hence
the two formulations attain the same objective. There are $O(B)$ budget
layers per block, $O(\min(B,m_g))$ local budget choices, $O(m_g^2)$
endpoint choices, and $O(m_g)$ interior variables per endpoint edge,
giving the bound.
\end{proof}

The decomposition is constructive: sample outgoing edges proportionally
to their flows and sample a fixed-cardinality subset with the conditional
interior marginals. For the latter, cumulative marginal intervals and a
uniform offset in $[0,1)$ give systematic sampling with the required
cardinality and marginal probabilities. Preserve the original coordinate
indices when returning the queried complements. The implementation reports
floating-point feasibility and marginal repairs; its numerical solution
is not an outward-rounded proof certificate.

\subsection{A sharp separation from deterministic adaptive recovery}
\begin{proposition}[One-block randomization gap]
\label{prop:random-gap}
For one exact-total block with $E>0$, $m$ missing coordinates, and
$0\leq B\leq m-2$, put $r=m-B$ and
$h=\min_{|S|=r}d(S)$. Let $D_B=Eh$ be its optimal deterministic
worst-case width, which also equals the optimum over deterministic
adaptive policies. Then
\begin{equation}
 \frac{r}{m}D_B\leq R_B\leq D_B.
 \label{eq:random-gap}
\end{equation}
When $r$ divides $m$, the lower bound is attained by instances with
$D_B>0$, so the factor $m/r$ cannot be improved in general.
\end{proposition}
\begin{proof}
Pad every batch to $B$ queries. Each retains $r$ coordinates, and therefore
\[
 \sum_{i=1}^m\mathbb E_\pi v_i(Q)
 =r\mathbb E_\pi d(S(Q))\geq rh.
\]
Its largest expected coordinate coefficient is at least $rh/m$;
\eqref{eq:random-value} proves the lower bound. A point mass at the
deterministic optimal batch proves the upper bound. Equality of
deterministic batch and adaptive optima follows from the preceding
zero-answer argument.

For sharpness, form $m/r$ coefficient clusters, each containing $r$
equally spaced points of range $h>0$, and separate adjacent clusters
by gaps strictly greater than $h$. Every $r$-point subset has range
at least $h$, and a complete cluster attains it. Uniformly retain one
cluster and query its complement. Each coordinate then has expected
coefficient $hr/m$, giving $R_B=Ehr/m$.
\end{proof}

If $D_B=0$, both optima are zero and their ratio is undefined.
The separation concerns randomized nonadaptive batches versus all
deterministic adaptive policies; it does not establish optimality among
randomized adaptive policies. Every realized sharp comparison interval
contains the comparison value for the fixed record and for every completion
consistent with its query answers. The improvement is in worst expected width, not a bound
on every realized width, a higher certification probability, or measured
retrieval savings. The result does not extend here to overlapping or
uncertain totals, physical caps, or heterogeneous query costs.
